%-----------------------------------
% Meta Informationen zur Arbeit
%-----------------------------------
% Diese Vorlage eignet sich fuer verschiedene Arten wissenschaftlicher Arbeiten
% (Bachelor-/Master-Thesis, Hausarbeit, Seminararbeit, Projektarbeit, Exposé ...).
% Die meisten Anpassungen erfolgen hier in dieser Datei.

% ifthen wird fuer die Fallunterscheidung auf der Titelseite benoetigt.
\usepackage{ifthen}

% Autor
\newcommand{\myAutor}{Max Mustermann}

% Adresse
\newcommand{\myAdresse}{Heidestra\ss e 17 \\ \> \> \> 51147 Köln}

% Titel der Arbeit
\newcommand{\myTitel}{LATEX-Vorlage - mit Biblatex}

%-----------------------------------
% Art der Arbeit (Modus-Schalter)
%-----------------------------------
% \myArbeitsModus steuert automatisch das Layout der Titelseite:
%   thesis            -> Abschlussarbeit: zeigt den angestrebten akademischen
%                        Grad sowie Erst-/Zweitgutachter
%                        (z. B. Bachelor-/Master-Thesis).
%   lehrveranstaltung -> Arbeit im Rahmen einer Lehrveranstaltung: zeigt die
%                        Lehrveranstaltung sowie den Betreuer
%                        (z. B. Hausarbeit, Seminararbeit, Projektarbeit, Exposé).
\newcommand{\myArbeitsModus}{thesis}

% Bezeichnung der Arbeit, wie sie auf der Titelseite erscheint.
% Beispiele: Bachelor Thesis, Master Thesis, Hausarbeit, Seminararbeit,
%            Projektarbeit, Exposé zur Bachelorarbeit
\newcommand{\myThesisArt}{Bachelor Thesis}

% Betreuer bzw. Erstgutachter
% (die Beschriftung auf der Titelseite richtet sich nach \myArbeitsModus)
\newcommand{\myBetreuer}{Prof. Dr. Peter Lustig}

% Zweitgutachter (nur im Modus 'thesis' relevant).
% Leer lassen ({}), wenn kein Zweitgutachter auf der Titelseite erscheinen soll.
\newcommand{\myZweitgutachter}{}

% Lehrveranstaltung (nur im Modus 'lehrveranstaltung' relevant)
\newcommand{\myLehrveranstaltung}{Modul Nr. 1}

% Zu erlangender akademischer Grad (nur im Modus 'thesis' relevant)
\newcommand{\myAkademischerGrad}{Bachelor of Science (B.Sc.)}

% Matrikelnummer
\newcommand{\myMatrikelNr}{123456}

% Ort
\newcommand{\myOrt}{Düsseldorf}

% Datum der Abgabe
\newcommand{\myAbgabeDatum}{\today}

% Semesterzahl
\newcommand{\mySemesterZahl}{7}

% Name der Hochschule
\newcommand{\myHochschulName}{FOM Hochschule für Oekonomie \& Management}

% Standort der Hochschule
\newcommand{\myHochschulStandort}{Düsseldorf}

% Studiengang
\newcommand{\myStudiengang}{Wirtschaftsinformatik}

% Firma
\newcommand{\myFirma}{Mustermann GmbH}

% Definition der Sprache
\ifdefined\FOMEN
%Englisch
\entrue
\usepackage[english]{babel}
\else
%Deutsch
\detrue
\usepackage[ngerman]{babel}
\fi
